\documentclass{article}
\usepackage[margin=1in]{geometry}

\title{ECE6255 - Team 9 Report}
\author{Jack Newcomb, Anish Nair, Hao Ma}
\date{\today}

\begin{document}

\maketitle

\section{Introduction}

The goal of this project is to develop a software system for modifying the duration of selected segments of a speech signal. The framework is designed to support user-specified edits in either of two forms: a scaling factor applied to a selected segment, or a desired target duration for that segment. In addition, the system supports multiple edits within a single audio file, and is intended to perform duration modification while preserving perceptually important speech characteristics as well as possible.

To support this goal, a modular processing pipeline was designed that loads an input audio signal, parses one or more requested edits, applies duration modification to the selected segments, and writes the resulting output signal. This structure separates the signal-processing algorithm from the surrounding application logic, making the system easier to test, validate, and extend.

For the core duration-modification stage, we selected the Pitch-Synchronous Overlap and Add (PSOLA) method. PSOLA is a well-known speech processing technique for time-scale modification because it allows changes in segment duration while better preserving pitch-related structure than naive resampling approaches. In this project, PSOLA serves as the primary algorithmic method for modifying speech duration within the broader editing framework.

\section{System Overview}

This project was implemented in Python using a modular design to separate responsibilities across the major stages of the processing pipeline. At a high level, the system takes an input speech waveform, accepts one or more user-specified duration edits, applies duration modification to the selected segments, and reconstructs a final edited output waveform. This organization was chosen to keep the signal-processing algorithm isolated from file handling, edit parsing, and reconstruction logic, making the framework easier to test and extend.

The system is composed of the following primary modules:

\begin{itemize}
  \item \textbf{Audio I/O:} Responsible for loading input WAV files and writing processed output audio files.
  
  \item \textbf{Edit Specification Interface:} Handles the representation and parsing of user-requested edits, including start time, end time, and either a scaling factor or target duration.
  
  \item \textbf{Segment Extraction:} Converts time-based edit definitions into sample indices and extracts the corresponding waveform segments to be modified.
  
  \item \textbf{Duration-Scaling Interface:} Provides the interface between the general editing framework and the core duration-modification algorithm. In the final system, this module is responsible for applying PSOLA-based time-scale modification to the selected speech segment.
  
  \item \textbf{Stitching/Reconstruction:} Reinserts the modified segment into the original waveform while preserving the unedited regions of the signal.
  
  \item \textbf{Output Generation:} Produces the final edited waveform and any associated artifacts needed for inspection or analysis.
\end{itemize}

In addition to the core pipeline, supporting scripts were developed for post-processing and evaluation. These scripts assist with generating analysis outputs such as waveform and spectrogram comparisons, and help support the final assessment of system performance.

This modular structure made it possible to validate much of the end-to-end workflow independently of the final algorithm implementation. As a result, components such as file handling, edit validation, batch processing, and output generation could be developed and tested before the PSOLA stage was fully integrated.

\section{Software Architecture and Design}

This section provides a more detailed discussion of the software engineering decisions behind the implementation of each module. In addition to the underlying speech-processing objective, the project required a software structure that could support user-configurable edits, robust validation, modular testing, and straightforward integration of the core duration-modification algorithm. For this reason, the system was designed as a pipeline of relatively isolated modules with clearly defined responsibilities and interfaces.

At the application level, a main orchestration layer coordinates the flow between modules. It loads the input waveform, parses and validates requested edits, extracts the relevant segment, invokes the duration-scaling interface, reconstructs the waveform, and finally writes the edited output. This orchestration layer keeps module interactions explicit while avoiding unnecessary entanglement between user-facing control flow and lower-level processing functions.

\subsection{Audio I/O}

Audio input and output functionality was implemented in \texttt{audio\_io.py}. This module is responsible for reading input \texttt{.wav} files into an internal array representation and writing the final processed waveform back to disk. Separating file handling from the remainder of the processing pipeline simplified both debugging and testing, since waveform manipulation could be developed independently of the details of file parsing.

The \texttt{load\_wav} function loads a \texttt{.wav} file and performs extensive validation to guard against invalid inputs and common edge cases. These checks include file existence, invalid or unsupported file contents, empty audio data, and invalid sample rate conditions. By concentrating such validation in a single module, the remainder of the pipeline can operate under clearer assumptions about the structure and validity of the loaded signal.

The \texttt{save\_wav} function handles writing the processed waveform to a \texttt{.wav} file. This provides a clean endpoint for the editing pipeline and ensures that the output stage is separated from the signal-processing logic itself.

\subsection{Edit Specification Interface}

The edit specification interface is responsible for representing and validating user-requested duration modifications before they are applied to the waveform. Conceptually, this layer serves as the boundary between the user-facing workflow and the lower-level signal-processing pipeline. Each requested edit specifies a time interval within the original signal together with either a duration scaling factor or a target output duration.

A key design goal of this interface was flexibility. The system was designed to support both single-edit and batch-edit workflows, allowing a user either to modify one segment directly or to provide multiple requested edits through a structured input format such as a CSV file. This made the framework more practical to use while preserving a consistent internal representation of edits once parsing was complete.

Validation was also an important responsibility of this module. Because downstream processing assumes well-formed edit requests, the edit interface performs checks such as interval consistency, nonnegative timing, and detection of overlapping edit regions in batch mode. This prevents invalid requests from propagating further into the processing pipeline and improves the overall robustness of the system.

\subsection{Segment Extraction}

Segment extraction was implemented in \texttt{segment.py}. This module is responsible for converting user-specified time intervals into the corresponding sample-index ranges and extracting the selected waveform segment from the full audio signal. The extracted segment can then be passed directly to the duration-modification module for processing.

This separation was useful because it isolates the time-to-sample conversion logic from both the user interface layer and the signal-processing algorithm. In doing so, it reduces coupling and makes the duration-modification module easier to test independently, since that module can operate on a clean one-dimensional sample array rather than needing to manage edit parsing or file-level concerns.

The segment extraction stage also includes several validation checks, including checks for negative times, invalid interval definitions, and empty extracted sample ranges. The output of this module is a one-dimensional array of samples together with the implicit location information needed for later reinsertion into the original waveform.

\subsection{Duration-Scaling Interface}

The duration-scaling interface is the point at which the general editing framework connects to the core speech-processing algorithm. Its purpose is to accept an extracted waveform segment together with the requested edit parameters, apply duration modification, and return a modified segment suitable for reintegration into the original signal.

This abstraction was intentionally designed so that the surrounding application logic would remain independent of the specific implementation details of the underlying algorithm. In the final system, this module is responsible for applying the PSOLA-based duration-modification method. By isolating the algorithm behind a clear interface, the remainder of the system could be developed and tested before the final PSOLA implementation was completed.

\subsection{Stitching/Reconstruction}

Stitching and reconstruction are performed after the duration-modification stage returns a modified waveform segment. This functionality was implemented in \texttt{stitch.py}. The purpose of this module is to replace the original user-specified segment with the modified segment while preserving the portions of the waveform that were not edited.

This design keeps reconstruction logic separate from both extraction and duration modification. As a result, each stage of the pipeline performs a narrow and well-defined task: extraction isolates the region of interest, the algorithm modifies that region, and reconstruction reinserts the result into the full signal. This modularity improves readability and makes it easier to verify correctness at each stage.

\subsection{Output Generation}

Final output generation consists primarily of producing the edited waveform and associated analysis artifacts. The edited audio output is generated through the same audio I/O layer used for file writing, while additional post-processing scripts were developed for evaluation and inspection of results.

In particular, supporting scripts were used to generate artifacts such as waveform comparisons, spectrogram comparisons, and other plots relevant to qualitative and quantitative assessment. Although these evaluation tools are not part of the core editing pipeline itself, they form an important part of the overall project because they provide a structured way to assess the effectiveness of the implemented duration-modification method.

\section{Configuration and Usage}

The project was designed to support straightforward command-line use for both single-edit and batch-edit duration modification tasks. The implementation assumes a mono \texttt{.wav} input file and accepts user-specified edit parameters that determine which portion of the waveform should be modified and how the duration change should be applied.

\subsection{Environment and Dependencies}

The recommended execution environment for the project is Python 3.10 or later. The primary dependencies used by the framework are \texttt{numpy}, \texttt{scipy}, \texttt{matplotlib}, and \texttt{pytest}. These libraries support waveform manipulation, file I/O, visualization, and testing.

Dependencies may be installed with:

\begin{verbatim}
pip install numpy scipy matplotlib pytest
\end{verbatim}

\subsection{Input Assumptions}

The system operates on mono \texttt{.wav} speech files. Each requested edit is specified by a start time \(t_1\) and end time \(t_2\), together with either a duration scaling factor or a target output duration. Internally, these time-domain edit requests are converted into sample-index ranges before segment extraction and processing.

The framework supports two modes of operation:
\begin{itemize}
    \item \textbf{Single-edit mode}, in which one speech segment is modified at a time.
    \item \textbf{Batch-edit mode}, in which multiple edits are provided through a CSV specification file.
\end{itemize}

\subsection{Single-Edit Usage}

Single-edit mode is used when modifying one selected region of the input waveform. In this mode, the user specifies the input file, output file, time interval, and either a scale factor or a target duration.

A scale-factor-based edit may be invoked as follows:

\begin{verbatim}
python main.py examples/tsignal.wav --output examples/output.wav \
    --t1 0.20 --t2 0.40 --scale 1.5
\end{verbatim}

This command stretches the region from \(0.20\) s to \(0.40\) s by a factor of \(1.5\).

Alternatively, a target duration may be specified directly:

\begin{verbatim}
python main.py examples/tsignal.wav --output examples/output.wav \
    --t1 0.20 --t2 0.40 --target-duration 0.50
\end{verbatim}

In this case, the selected segment is modified so that its output duration is \(0.50\) s.

\subsection{Batch-Edit Usage}

Batch-edit mode allows multiple edits to be applied to the same input file using a CSV specification. This is useful when several independent regions of the waveform must be modified in a single run.

Batch mode may be invoked as follows:

\begin{verbatim}
python main.py examples/tsignal.wav --output examples/output.wav \
    --csv examples/sample_input.csv
\end{verbatim}

The CSV file must include the columns \texttt{t1} and \texttt{t2}, along with at least one of \texttt{scale} or \texttt{target\_duration}. Each row must specify exactly one of these two duration-modification options.

An example CSV file is:

\begin{verbatim}
t1,t2,scale,target_duration
0.20,0.40,1.5,
0.60,0.80,,0.30
\end{verbatim}

In this example, the first segment is stretched by a factor of \(1.5\), while the second segment is modified so that its new duration is \(0.30\) s.

\subsection{Batch Timing Convention and Validation}

A key design choice in batch mode is that edit times are interpreted relative to the original input-file timeline rather than the progressively modified output timeline. This means that if an earlier edit changes the duration of one region, later edit times are still understood as referring to their intended locations in the original signal. Internally, the system remaps these later edit boundaries onto the current modified waveform as needed.

To preserve correctness and avoid ambiguity, overlapping edits in the original timeline are not supported. For example, the following CSV specification is invalid:

\begin{verbatim}
t1,t2,scale
0.20,0.50,1.2
0.40,0.70,0.8
\end{verbatim}

because the second edit overlaps the first in the original signal.

\subsection{Evaluation and Testing Workflow}

In addition to the main processing pipeline, the repository includes a separate evaluation script for comparing original and modified audio outputs. Basic evaluation may be run with:

\begin{verbatim}
python scripts/evaluate.py --original examples/tsignal.wav \
    --modified examples/output.wav --report-dir reports/basic_eval
\end{verbatim}

This generates waveform comparisons, spectrogram visualizations, and a basic metrics report. Region-focused evaluation is also supported by specifying the edited interval and a margin around it.

The project also includes an automated test suite, which may be run with:

\begin{verbatim}
pytest
\end{verbatim}

These tests cover key framework functionality, including segment extraction, stitching behavior, CSV parsing, batch processing, target-duration handling, remapping of batch edit timelines, and rejection of overlapping edits.







\section{Duration Modification Algorithm}

\section{Evaluation Methodology}

\section{Results}

\section{Discussion and Limitations}

\section{Conclusion}

This project developed a modular software framework for arbitrary duration modification of selected segments within a speech waveform. The system was designed to support flexible user-specified edits through either scale-factor-based or target-duration-based input, while maintaining a clear separation between waveform processing, edit specification, segment extraction, reconstruction, and evaluation.

A central goal of the implementation was to build an end-to-end pipeline that could support both algorithm development and practical use. To that end, the project provides functionality for mono WAV input/output, single-edit and batch-edit workflows, validation of user inputs, rejection of overlapping edit regions, and supporting tools for evaluation and testing. This modular design allowed substantial portions of the framework to be implemented and verified independently of the final duration-modification algorithm.

The selected core method for duration modification is PSOLA, a well-established speech processing technique for time-scale modification of voiced speech. Within the broader software framework, PSOLA serves as the principal algorithmic mechanism for changing segment duration while seeking to preserve perceptually important speech characteristics. The surrounding system was intentionally structured to make this algorithm easier to integrate, test, and evaluate.

Overall, the project demonstrates both the signal-processing and software-engineering components required for arbitrary segment duration modification in speech. In addition to addressing the immediate project goals, the resulting framework provides a foundation that can be extended in future work to support more advanced duration-modification methods, additional speech characteristics, and more sophisticated evaluation procedures.

\end{document}